\documentclass{article}
\usepackage{amsmath,amssymb,amsthm}
\usepackage{xcolor}
\usepackage{ctex}

\newtheorem{proposition}{命题}
\newtheorem{corollary}{推论}

\title{为什么更大的梯度范数意味着更小的 $\sigma$}
\author{}
\date{}

\begin{document}
\maketitle

\section{背景}

考虑带有噪声函数评估的 antithetic ES 梯度估计器，对于 $f:\mathbb{R}^d\to\mathbb{R}$，我们只能观测到 $F(x) = f(x) + \xi$，其中 $\mathbb{E}[\xi\mid x]=0$ 且 $\mathbb{E}[\xi^2\mid x]\leq\gamma^2$：
\[
\widehat{g}_\sigma(x) = \frac{1}{m}\sum_{i=1}^m \frac{F(x+\sigma u_i) - F(x-\sigma u_i)}{2\sigma}u_i,
\quad u_i\sim\mathcal{N}(0,I_d).
\]

均方误差的上界为：
\[
\mathbb{E}\|\widehat{g}_\sigma(x) - \nabla f(x)\|^2
\leq
L^2\sigma^2 d + \frac{d}{m}\|\nabla f(x)\|^2 + \frac{L^2\sigma^2 d^3}{m} + \frac{\gamma^2 d}{m\sigma^2}.
\]

\section{主要结果}

\subsection{最小化绝对误差：$\sigma$ 与 $\|\nabla f\|$ 无关}

\begin{proposition}[绝对误差最小化]
最小化 $\mathbb{E}\|\widehat{g}_\sigma - \nabla f\|^2$ 的最优 $\sigma$ 为
\[
\sigma^* \propto \left(\frac{\gamma^2}{mL^2}\right)^{1/4},
\]
它与 $\|\nabla f(x)\|$ \textcolor{red}{无关}。
\end{proposition}

\begin{proof}
在四项误差中，只有第一项和最后一项依赖于 $\sigma$。平衡主导项：
\[
L^2\sigma^2 d \approx \frac{\gamma^2 d}{m\sigma^2}
\implies
L^2\sigma^4 \approx \frac{\gamma^2}{m}
\implies
\sigma^* \propto \left(\frac{\gamma^2}{mL^2}\right)^{1/4}.
\]
注意 $\|\nabla f(x)\|$ 只出现在第二项 $\frac{d}{m}\|\nabla f\|^2$ 中，而该项不依赖于 $\sigma$。
\end{proof}

\subsection{保持恒定信噪比：$\sigma \propto 1/\|\nabla f\|$}

\begin{proposition}[信噪比视角]
如果我们希望在整个训练过程中保持恒定的信噪比，那么
\[
\boxed{\sigma \propto \frac{1}{\|\nabla f(x)\|}},
\]
即，\textcolor{red}{梯度范数越大需要越小的 $\sigma$}。
\end{proposition}

\begin{proof}
信号强度为 $\|\mathbb{E}[\widehat{g}_\sigma]\| \approx \|\nabla f(x)\|$（忽略 $O(\sigma^2)$ 偏差）。

噪声方差（忽略高阶项）为：
\[
\text{Var}(\widehat{g}_\sigma) \approx \frac{d}{m}\|\nabla f\|^2 + \frac{\gamma^2 d}{m\sigma^2}.
\]

定义信噪比：
\[
\text{SNR}^2 = \frac{\text{Signal}^2}{\text{Noise}} = \frac{\|\nabla f\|^2}{\frac{d}{m}\|\nabla f\|^2 + \frac{\gamma^2 d}{m\sigma^2}}.
\]

为了保持 $\text{SNR} = C$（常数），我们要求：
\[
\|\nabla f\|^2 = C^2\left[\frac{d}{m}\|\nabla f\|^2 + \frac{\gamma^2 d}{m\sigma^2}\right].
\]

重新整理：
\[
\|\nabla f\|^2\left(1 - \frac{C^2 d}{m}\right) = \frac{C^2\gamma^2 d}{m\sigma^2}.
\]

假设 $C^2 < m/d$（SNR 不能无限大），求解 $\sigma$：
\[
\sigma^2 = \frac{C^2\gamma^2 d}{m\|\nabla f\|^2\left(1 - \frac{C^2 d}{m}\right)}
\propto \frac{1}{\|\nabla f\|^2}.
\]

因此：
\[
\boxed{\sigma \propto \frac{1}{\|\nabla f\|}.}
\]
\end{proof}

\subsection{噪声放大分析：$\sigma \propto 1/\|\nabla f\|$}

\begin{proposition}[噪声放大]
从有限差分的视角，单方向信噪比为
\[
\text{SNR}_{\text{single}} = \frac{\|\nabla f\|\cdot\sigma}{\gamma}.
\]
为了保持恒定的 $\text{SNR}_{\text{single}}$，我们需要 $\sigma \propto 1/\|\nabla f\|$。
\end{proposition}

\begin{proof}
沿方向 $u$ 的有限差分近似为：
\[
\frac{F(x+\sigma u) - F(x-\sigma u)}{2\sigma} = \frac{f(x+\sigma u) - f(x-\sigma u)}{2\sigma} + \frac{\xi_+ - \xi_-}{2\sigma}.
\]

信号部分：
\[
\frac{f(x+\sigma u) - f(x-\sigma u)}{2\sigma} \approx \nabla f(x) \cdot u,
\]
其大小为 $\|\nabla f\|$（在 $u$ 上的期望）。

噪声部分：
\[
\frac{\xi_+ - \xi_-}{2\sigma}
\]
的标准差为：
\[
\frac{\sqrt{\text{Var}(\xi_+) + \text{Var}(\xi_-)}}{2\sigma} = \frac{\gamma\sqrt{2}}{2\sigma} \approx \frac{\gamma}{\sigma}.
\]

因此，单方向 SNR 为：
\[
\text{SNR}_{\text{single}} = \frac{\|\nabla f\|}{\gamma/\sigma} = \frac{\|\nabla f\|\cdot\sigma}{\gamma}.
\]

为了保持 $\text{SNR}_{\text{single}} = C$（常数）：
\[
\|\nabla f\|\cdot\sigma = C\gamma
\implies
\boxed{\sigma = \frac{C\gamma}{\|\nabla f\|}}.
\]

这表明 \textcolor{red}{$\sigma$ 应该随 $\|\nabla f\|$ 增大而\emph{减小}}。
\end{proof}

\subsection{一步优化误差：$\sigma \propto 1/\|\nabla f\|^{1/2}$}

\begin{proposition}[一步更新误差]
如果我们执行一步梯度下降 $x' = x - \alpha\widehat{g}_\sigma(x)$ 并最小化期望的一步误差，最优 $\sigma$ 满足：
\[
\sigma \propto \frac{1}{\|\nabla f\|^{1/2}}.
\]
\end{proposition}

\begin{proof}
考虑使用学习率 $\alpha$ 的单步更新。实际更新与理想更新之间的差异可由以下式子刻画：
\[
\mathbb{E}[f(x')] - f(x - \alpha\nabla f)
\approx
\alpha^2\sigma^2\|\nabla f\|^2 + \frac{\alpha^2 L}{2}\left[\frac{d}{m}\|\nabla f\|^2 + \frac{\gamma^2 d}{m\sigma^2}\right].
\]

依赖于 $\sigma$ 的项为：
\[
\text{Error}(\sigma) = \alpha^2\sigma^2\|\nabla f\|^2 + \frac{\alpha^2 L\gamma^2 d}{2m\sigma^2}.
\]

对 $\sigma$ 求导并令其为零：
\[
\frac{\partial}{\partial\sigma}\left[\alpha^2\sigma^2\|\nabla f\|^2 + \frac{\alpha^2 L\gamma^2 d}{2m\sigma^2}\right] = 0,
\]
\[
2\alpha^2\sigma\|\nabla f\|^2 - \frac{\alpha^2 L\gamma^2 d}{m\sigma^3} = 0,
\]
\[
\sigma^4 = \frac{L\gamma^2 d}{2m\|\nabla f\|^2}.
\]

因此：
\[
\boxed{\sigma \propto \frac{1}{\|\nabla f\|^{1/2}}.}
\]

再一次，\textcolor{red}{$\sigma$ 应该随 $\|\nabla f\|$ 增大而减小}。
\end{proof}

\section{总结}

\begin{corollary}[反直觉的结果]
在加性噪声模型 $F(x) = f(x) + \xi$ 下，多个理论视角表明：
\begin{itemize}
\item \textbf{绝对误差最小化}：$\sigma$ 与 $\|\nabla f\|$ 无关。
\item \textbf{恒定信噪比}：$\sigma \propto 1/\|\nabla f\|$。
\item \textbf{噪声放大}：$\sigma \propto 1/\|\nabla f\|$。
\item \textbf{一步误差}：$\sigma \propto 1/\|\nabla f\|^{1/2}$。
\end{itemize}

\textcolor{red}{\textbf{这些结果都不支持"$\|\nabla f\|$ 越大需要越大的 $\sigma$"这一猜想。}}

事实上，当梯度范数很大时（信号强），我们可以使用\emph{更小}的 $\sigma$ 来获得更精确的梯度估计，而不会被噪声主导。
\end{corollary}

\section{直觉解释}

关键洞察是：
\[
\text{SNR} \propto \|\nabla f\| \cdot \sigma.
\]

当 $\|\nabla f\|$ 增大时：
\begin{itemize}
\item 信号变得更强。
\item 即使 $\sigma$ 固定不变，SNR 也会提高。
\item 我们可以\emph{减小} $\sigma$ 来获得更精确的梯度估计（减小有限差分的偏差），同时仍然保持可接受的 SNR。
\end{itemize}

反过来，当 $\|\nabla f\|$ 很小时（接近收敛）：
\begin{itemize}
\item 信号很弱。
\item 噪声主导：$\frac{\xi_+ - \xi_-}{2\sigma}$ 中的 $1/\sigma$ 放大因子变得至关重要。
\item 我们应该增大 $\sigma$，或增大采样方向数 $m$，或停止训练。
\end{itemize}

\end{document}
